\chapter{信息访问与条件积分估计}
\label{app:estimation}
\section{四个积分对象与可观测信息}
\begin{proposition}[双向 KL 的数值与梯度积分]
有限词表上，$P=P_\theta>0$ 可微且归一化，$Q>0$ 冻结。令 $\ell=\log(P/Q)$、$\psi=\nabla_\theta\log P$，则
\begin{equation}
\KL(P\Vert Q)=\E_P\ell,\quad \nabla_\theta\KL(P\Vert Q)=\E_P(\ell\psi),\quad
\nabla_\theta\KL(Q\Vert P)=-\E_Q\psi.
\end{equation}
\end{proposition}
\begin{proof}
本附录局部记当前学生为 $P=P_\theta$、冻结参考为 $Q$、动作提案为 $\beta$；三者职责不同。令 $\ell_a=\log(P_a/Q_a)$、$\psi_a=\nabla_\theta\log P_a$。双向 KL 及其梯度为
\begin{equation}
K=\E_P\ell,\quad g=\E_P[\ell\psi],\qquad
K^{\rm F}=-\E_Q\ell,\quad g^{\rm F}=-\E_Q\psi.
\label{eq:four-integrals}
\end{equation}
$g$ 的求导先产生 $\E_P[(\ell+1)\psi]$，再利用 $\E_P\psi=0$。参考为锚点时可代入 $Q=q$；BPO 学生中心 $\KL(q\Vert p)$ 对应这里的正向积分。
\end{proof}

完整词表概率允许精确有限求和。头部查询只给出集合 $S$ 上的 $Q_a$ 和尾部总质量；实际动作查询再给出该动作的一项概率。若 $S$ 是教师真正的 Top-$K_{\rm head}$，未知尾部还满足排名约束 $Q_a\le\min_{b\in S}Q_b$。接口是否确定某个量，应检查它在全部符合这些观测的分布上是否相同。

只保留桶质量时，不能恢复桶内形状。若实际动作落在尾部，其额外概率观测必须计入信息；也不能把一次尾部查询当成全部尾部已知。固定支持集的微分与每步重新选 Top-K 的程序是不同对象，以下推导均冻结分区。

\section{粗粒化、下界与不可识别性}
\begin{theorem}[分区链式分解与细化单调性]
设 $\mathcal C$ 为固定词表分区，$P,Q>0$，桶质量与条件分布分别为 $P_C,Q_C,P^C,Q^C$。则
\begin{equation}
\KL(P\Vert Q)=\sum_CP_C\log(P_C/Q_C)+\sum_CP_C\KL(P^C\Vert Q^C).
\end{equation}
交换 $P,Q$ 给另一方向；细化分区使粗粒度 KL 非减，完整 singleton 分区恢复完整 KL。
\end{theorem}
\begin{proof}
设 $\mathcal C$ 为词表分区，定义 $P_C=\sum_{a\in C}P_a$、$Q_C=\sum_{a\in C}Q_a$、$\ell_C=\log(P_C/Q_C)$、$P^C_a=P_a/P_C$、$Q^C_a=Q_a/Q_C$。逐点有
\begin{equation}
\ell_a=\ell_C+\log(P_a^C/Q_a^C).
\end{equation}
乘 $P_a$，先桶内求和再跨桶求和，得到
\begin{equation}
\KL(P\Vert Q)=K_{\mathcal C}+\sum_CP_C\KL(P^C\Vert Q^C),
\qquad K_{\mathcal C}=\sum_CP_C\ell_C.
\label{eq:partition-chain}
\end{equation}
交换 $P,Q$ 给出正向版本。细化分区时，新旧粗化 KL 的差为每个旧桶内细桶标签 KL 的加权和，故非负。全 singleton 分区恢复完整值。
\end{proof}

\begin{proposition}[桶质量信息下的 KL 下包络]
固定 $P>0$ 与各桶正质量 $Q_C$，仅施加桶质量约束时，$\inf_Q\KL(P\Vert Q)=\sum_CP_C\log(P_C/Q_C)$，唯一最优补全为 $Q_a=Q_CP_a^C$。交换 KL 方向仍由同一补全达到相应粗粒度下界。加入 Top-K 排名后，该下界不一定仍可达到。
\end{proposition}
\begin{proof}
在只知道 $Q_C$ 的信息类上，$K_{\mathcal C}$ 是完整 KL 的锐下包络，由 $Q_a=Q_CP_a^C$ 唯一达到；两个方向同时成立。加入真实 Top-K 排名后该构造可能不可行。例如 $P=(0.1,0.8,0.1)$，已知 $Q_1=0.4,Q_2+Q_3=0.6$ 且动作一是 Top-1；质量约束最优构造给 $Q_2=8/15>0.4$，违反排名。因此 Tail-aware 仍是下界，但不必是完整信息类上的锐下界。
\end{proof}

\begin{proposition}[稀疏头部信息不控制完整逆向 KL 上界]
当至少两个尾部概率仍未知时，存在具有相同真实 Top-1 概率和尾部质量的正参考分布族，使对同一个固定 $P$ 的 $\KL(P\Vert Q)$ 无上界。
\end{proposition}
\begin{proof}
相同头部信息也不控制逆向 KL 上界。固定 $P=(0.4,0.2,0.4)$ 与 $Q_\epsilon=(0.7,\epsilon,0.3-\epsilon)$、$0<\epsilon<0.15$，Top-1 和尾部质量完全相同，而 $0.2\log(0.2/\epsilon)\to\infty$。该反例依赖至少两个仍未知的尾部坐标。
\end{proof}

\begin{proposition}[粗粒度 KL 的梯度残差]
固定分区和参考 $Q$，令 $\ell_C=\log(P_C/Q_C)$、$\bar\psi_C=\E_{P^C}\psi$，则
\begin{align}
\nabla K_{\mathcal C}&=\sum_CP_C\ell_C\bar\psi_C,\\
g-\nabla K_{\mathcal C}
&=\sum_CP_C\{\Cov_{P^C}(\ell-\ell_C,\psi)+\KL(P^C\Vert Q^C)\bar\psi_C\}.
\end{align}
标量粗化误差的单调改进不一般推出梯度偏差范数的单调改进。
\end{proposition}
\begin{proof}
记 $\bar\psi_C=\nabla\log P_C=\E_{P^C}\psi$。对粗化目标求导：
\begin{align}
\nabla K_{\mathcal C}
&=\sum_C(\ell_C+1)\nabla P_C
=\sum_CP_C\ell_C\bar\psi_C,\\
g-\nabla K_{\mathcal C}
&=\sum_CP_C\{\Cov_{P^C}(\ell-\ell_C,\psi)
+\KL(P^C\Vert Q^C)\bar\psi_C\}.
\end{align}
最后一步对桶内乘积期望分解。标量误差只含条件 KL，梯度误差还依赖 score 相关性；标量下界改善不推出梯度偏差范数单调改善。
\end{proof}

\section{精确积分加残差抽样}
\begin{theorem}[统一估计与协方差]
设目标 $I_\nu(h)=\sum_a\nu_ah_a$，$h_a$ 可为向量。选择可精确积分的代理 $b_a$，令 $R_a=h_a-b_a$。若 $\nu_a\|R_a\|>0$ 时 $\beta_a>0$，则对 $A\sim\beta$，
\begin{equation}
\widehat I=I_\nu(b)+\frac{\nu_A}{\beta_A}R_A
\end{equation}
无偏，且
\begin{equation}
\Cov(\widehat I)=\sum_{\beta_a>0}\frac{\nu_a^2}{\beta_a}R_aR_a^\top
-I_\nu(R)I_\nu(R)^\top.
\label{eq:master-cov}
\end{equation}
\end{theorem}
\begin{proof}
随机修正的均值为 $\sum_{\beta_a>0}\nu_aR_a=I_\nu(R)$，支持条件使缺失动作贡献为零；与确定积分相加即得目标。随机项外积的期望为式~\eqref{eq:master-cov} 的求和项，减去均值外积得到协方差。取迹得到向量方差，无偏性使其等于均方误差。$M$ 个条件独立样本的平均协方差除以 $M$。
\end{proof}

\begin{proposition}[积分项、权重和支持的统一偏差]
记 $R=h-b$、$\mathcal D=\{a:\beta_a>0\}$。若实际估计器为 $I_\nu(b)+\delta+\widetilde w_AR_A$，则其偏差为
\begin{equation}
\delta+\sum_{\mathcal D}(\beta_a\widetilde w_a-\nu_a)R_a
-\sum_{\mathcal D^c}\nu_aR_a.
\end{equation}
\end{proposition}
\begin{proof}
若确定积分有误差 $\delta$，使用错误权重 $\widetilde w$ 或缺失支持 $\mathcal D^c$，偏差恰为
\begin{equation}
\delta+\sum_{a\in\mathcal D}(\beta_a\widetilde w_a-\nu_a)R_a
-\sum_{a\notin\mathcal D}\nu_aR_a.
\end{equation}
这是对实际期望与目标直接相减的结果。精确覆盖的零残差动作可以没有采样支持；未知残差则不能由有限权重恢复。
\end{proof}

\begin{proposition}[固定残差下的迹方差最优提案]
固定代理与残差 $R$，若 $L_R=\sum_a\nu_a\|R_a\|>0$，则最优提案为 $\beta_a^\star=\nu_a\|R_a\|/L_R$，最小迹方差为 $L_R^2-\|I_\nu(R)\|^2$。若 $L_R=0$，目标可确定性精确返回。
\end{proposition}
\begin{proof}
固定代理时，迹方差最优提案为
\begin{equation}
\beta_a^\star=\frac{\nu_a\|R_a\|}{\sum_b\nu_b\|R_b\|},\quad
V_{\min}=\left(\sum_a\nu_a\|R_a\|\right)^2-\|I_\nu(R)\|^2.
\end{equation}
证明对 $\sum_a[\nu_a\|R_a\|/\sqrt{\beta_a}]\sqrt{\beta_a}$ 用 Cauchy--Schwarz；取比例提案时等号成立。全部残差为零时直接返回确定积分。这个 oracle 提案通常需要未知尾部值，不能仅凭公式称其可实现。
\end{proof}

\begin{theorem}[多控制的 Gram 最优解]
设 $U_j$ 为二阶可积零均值控制，估计器为 $X-\sum_jc_jU_j$。令 $\Gamma_{ij}=\E\langle U_i,U_j\rangle$、$d_j=\E\langle U_j,X-\E X\rangle$。则全部迹方差最优系数为 $\Gamma^\dagger d+\ker\Gamma$，相对最优值的方差代价为 $(c-\Gamma^\dagger d)^\top\Gamma(c-\Gamma^\dagger d)$。
\end{theorem}
\begin{proof}
若有多个已知积分控制 $h_j$，令 $X=(\nu_A/\beta_A)h_A$、$U_j=(\nu_A/\beta_A)h_{j,A}-I_\nu(h_j)$、$\Gamma_{ij}=\E\langle U_i,U_j\rangle$、$d_j=\E\langle U_j,X-\E X\rangle$。估计器 $X-\sum_jc_jU_j$ 的方差为
\begin{equation}
V_0-2c^\top d+c^\top\Gamma c.
\end{equation}
若 $z\in\ker\Gamma$，则 $\sum_jz_jU_j=0$ 几乎处处，所以 $z^\top d=0$。因此 $d\in\operatorname{range}\Gamma$，最优系数为 $\Gamma^\dagger d+\ker\Gamma$，超额方差为 $(c-\Gamma^\dagger d)^\top\Gamma(c-\Gamma^\dagger d)$。这最小化协方差迹，不一般同时最小化每个方向。
\end{proof}

\section{头尾估计器作为统一定理的实例}
\begin{proposition}[三种确定性头部目标及其偏差]
固定非空头尾分区，令 $H=\sum_SP_a\ell_a$、$l=\log(\tau_P/\tau_Q)$、$B=\tau_P\KL(P^{\rm tail}\Vert Q^{\rm tail})$。完整逆向 KL 为 $H+\tau_Pl+B$。Truncated 与 Tail-aware 的偏差分别为 $-\tau_Pl-B,-B$；Normalized 优化头部条件 KL $H/P_S+\log(Q_S/P_S)$。
\end{proposition}
\begin{proof}
令 $\tau_P=P(S^c)$、$\tau_Q=Q(S^c)$、$J_A=\mathbf1\{A\notin S\}$，并定义
\begin{equation}
H=\sum_{a\in S}P_a\ell_a,\quad
l=\log(\tau_P/\tau_Q),\quad
B=\tau_P\KL(P^{\rm tail}\Vert Q^{\rm tail}).
\end{equation}
则 $K=H+\tau_Pl+B$。Truncated 返回 $H$，偏差为 $-\tau_Pl-B$；Tail-aware 返回 $H+\tau_Pl$，偏差为 $-B$；Normalized 返回
\begin{equation}
K_{\rm norm}=H/P_S+\log(Q_S/P_S),
\end{equation}
它是头部条件分布的 KL。后者改变积分权重，不能当成只删除尾部项。
\end{proof}

\begin{proposition}[标量尾部控制族的精确方差]
令 $A\sim P$，控制 $c$ 在抽样前固定。估计器 $\widehat K(c)=H+\tau_Pc+J_A(\ell_A-c)$ 无偏，且
\begin{equation}
\Var\widehat K(c)=\tau_P\Var_{\rm tail}\ell+P_S\tau_P(\E_{\rm tail}\ell-c)^2.
\end{equation}
因此最优常数为尾部条件均值；HETS 与桶残差校正分别取 $c=0,l$。
\end{proposition}
\begin{proof}
在 $A\sim P$ 下，标量控制族
\begin{equation}
\widehat K(c)=H+\tau_Pc+J_A(\ell_A-c)
\end{equation}
无偏。设 $\mu_\ell=\E_{P^{\rm tail}}\ell$、$V_\ell=\Var_{P^{\rm tail}}\ell$，全方差公式给出
\begin{equation}
\Var\widehat K(c)=\tau_PV_\ell+P_S\tau_P(\mu_\ell-c)^2.
\end{equation}
头部精确、尾部采样（Head-Exact Tail-Sampled，HETS）取 $c=0$，桶残差修正取 $c=l$，oracle 取 $c=\mu_\ell$。残差修正优于 HETS 当且仅当 $|\mu_\ell-l|<|\mu_\ell|$。
\end{proof}

\begin{proposition}[向量尾部控制族的协方差]
令 $h_g=\sum_SP_a\ell_a\psi_a$、$\mu_g=\E_{\rm tail}(\ell\psi)$，固定向量 $b$。则 $\widehat g(b)=h_g+\tau_Pb+J_A(\ell_A\psi_A-b)$ 对完整梯度无偏，协方差为
\begin{equation}
\tau_P\Cov_{\rm tail}(\ell\psi)+P_S\tau_P(\mu_g-b)(\mu_g-b)^\top.
\end{equation}
\end{proposition}
\begin{proof}
对梯度令 $h_g=\sum_SP_a\ell_a\psi_a$、$\mu_g=\E_{P^{\rm tail}}(\ell\psi)$、$\bar\psi_{\rm tail}=\nabla\log\tau_P$、$b_{\rm TA}=l\bar\psi_{\rm tail}$。统一向量控制为
\begin{equation}
\widehat g(b)=h_g+\tau_Pb+J_A(\ell_A\psi_A-b),
\end{equation}
\begin{equation}
\Cov\widehat g(b)=\tau_P\Cov_{\rm tail}(\ell\psi)
+P_S\tau_P(\mu_g-b)(\mu_g-b)^\top.
\end{equation}
HETS 为 $b=0$，ResOPD 为 $b=b_{\rm TA}$。常向量只改变头尾事件间的均值差，不改变尾部内部协方差。
\end{proof}

\begin{proposition}[正交双控制的最优值与排序]
on-policy 抽样下，取 $\chi=l(\psi-\bar\psi_{\rm tail})$。估计器 $\widehat g(b,\vartheta)=h_g+\tau_Pb+J_A(\ell_A\psi_A-b-\vartheta\chi_A)$ 无偏。若 $V_\chi>0$，最优值为 $b=\mu_g,\vartheta=D_\chi/V_\chi$，超额方差为
\begin{equation}
P_S\tau_P\|b-\mu_g\|^2+\tau_PV_\chi(\vartheta-D_\chi/V_\chi)^2.
\end{equation}
其中 $V_\chi=\E_{\rm tail}\|\chi\|^2$、$D_\chi=\E_{\rm tail}\langle\ell\psi-\mu_g,\chi\rangle$。SC-TA 优于 ResOPD 当且仅当 $2D_\chi>V_\chi$。
\end{proposition}
\begin{proof}
再加入尾部条件零均值方向 $\chi_a=l(\psi_a-\bar\psi_{\rm tail})$，定义
\begin{equation}
\widehat g(b,\vartheta)=h_g+\tau_Pb+J_A(\ell_A\psi_A-b-\vartheta\chi_A).
\end{equation}
尾部内部与事件控制正交，因为 $\E_{\rm tail}\chi=0$。因此最优 $b=\mu_g$，最优 $\vartheta=D_\chi/V_\chi$，其中
\begin{equation}
V_\chi=\E_{\rm tail}\|\chi\|^2,\qquad
D_\chi=\E_{\rm tail}\langle\ell\psi-\mu_g,\chi\rangle.
\end{equation}
配方给出精确代价
\begin{equation}
V(b,\vartheta)-V_{\min}
=P_S\tau_P\|b-\mu_g\|^2
+\tau_PV_\chi(\vartheta-D_\chi/V_\chi)^2.
\end{equation}
$V_\chi=0$ 时内部控制不起作用。SC-TA 对应 $b=b_{\rm TA},\vartheta=1$；它比 ResOPD 方差小当且仅当 $2D_\chi>V_\chi$。
\end{proof}

\begin{proposition}[精确头部不保证逐输入降方差]
对向量 integrand $u$、$I=\E_Pu$、$H_u=\sum_SP_au_a$，有
\begin{equation}
\Var(u_A)-\Var(H_u+J_Au_A)
=\sum_SP_a\|u_a\|^2-2\langle I,H_u\rangle+\|H_u\|^2.
\end{equation}
右侧可以为负；真实条件期望的 Rao--Blackwell 化则保证方差不增。
\end{proposition}
\begin{proof}
即使头部精确，方差也不必比直接抽样小。对任意 integrand $u$，令 $I=\E_Pu$、$H_u=\sum_SP_au_a$，直接展开二阶矩：
\begin{equation}
\Var(u_A)-\Var(H_u+J_Au_A)
=\sum_SP_a\|u_a\|^2-2\langle I,H_u\rangle+\|H_u\|^2.
\end{equation}
右侧可负。数值反例 $P=(0.18,0.819,0.001)$、$Q=(0.08,0.67,0.25)$、$S=\{2\}$（按一开始编号），直接 log-ratio 方差约为 $0.088913$，HETS 为 $0.129131$。真实 Rao--Blackwell 化应取条件期望 $\E_P[\ell\mid C]=\ell_C+\KL(P^C\Vert Q^C)$；只知道桶质量时缺少第二项。
\end{proof}

\section{提案变化、分层采样与正向 KL}
\begin{proposition}[一般提案下的最优尾部常控制]
提案 $\beta$ 覆盖尾部。令 $W=\sum_{S^c}P_a^2/\beta_a$、$U=\sum_{S^c}(P_a^2/\beta_a)h_a$、$T_h=\sum_{S^c}P_ah_a$。若 $W>\tau_P^2$，最优控制为 $(U-\tau_PT_h)/(W-\tau_P^2)$；若等号成立，$\beta=P^{\rm tail}$，控制逐样本抵消，$M$ 次独立抽样的方差为 $\tau_P^2\Var_{\rm tail}(h)/M$。
\end{proposition}
\begin{proof}
从 $\beta$ 采样时，所有残差改为 $J_A(P_A/\beta_A)(h_A-b)$。二阶矩权重是 $P^2/\beta$，所以原有 on-policy 最优控制通常改变。令
\begin{equation}
W=\sum_{S^c}P_a^2/\beta_a,\quad U=\sum_{S^c}(P_a^2/\beta_a)h_a,\quad
T_h=\sum_{S^c}P_ah_a,
\end{equation}
展开关于 $b$ 的二次式，最优 $b=(U-\tau_PT_h)/(W-\tau_P^2)$。若 $W=\tau_P^2$，Cauchy--Schwarz 等号要求 $\beta=P^{\rm tail}$，控制项逐样本抵消。此时直接尾部条件采样的方差为 $\tau_P^2\Var_{\rm tail}(h)/M$，不超过同样本数的 on-policy 头尾估计方差。
\end{proof}

\begin{proposition}[分层预算的最优分配]
各桶独立采样 $M_C$ 次，$\sigma_C^2=\Var_{P^C}h$。总估计方差为 $\sum_CP_C^2\sigma_C^2/M_C$。在连续预算 $\sum_CM_C=M$、$\sum_CP_C\sigma_C>0$ 下，最优分配为 $M_C\propto P_C\sigma_C$。控制若由数据估计，条件无偏结论还要求控制拟合与评估动作之间满足对应独立性。
\end{proposition}
\begin{proof}
更一般，分区内独立采样 $M_C$ 次，方差为 $\sum_CP_C^2\sigma_C^2/M_C$。固定连续预算 $\sum_CM_C=M$，Cauchy--Schwarz 给出最优分配 $M_C\propto P_C\sigma_C$；零方差桶直接返回常值，整数预算需另行取整。控制拟合与评估共用随机动作时，$\E U=0$ 不保证 $\E[c(A)U]=0$；应条件化于独立校准数据，或使用交叉拟合并保留折间相关性。
\end{proof}

\begin{proposition}[正向 KL 的头尾校正]
令 $H^{\rm F}=-\sum_SQ_a\ell_a$，则 $K^{\rm F}=H^{\rm F}-\tau_Ql+\tau_Q\KL(Q^{\rm tail}\Vert P^{\rm tail})$。若 $\beta$ 覆盖未知 $Q$ 尾部，则 $H^{\rm F}+\tau_Qc+J_A(Q_A/\beta_A)(-\ell_A-c)$ 无偏。粗粒度正向 KL 的真实梯度为 $-\sum_CQ_C\bar\psi_C$。
\end{proposition}
\begin{proof}
正向 KL 使用同一积分定理取 $\nu=Q,h=-\ell$ 或 $-\psi$。令 $H^{\rm F}=-\sum_SQ_a\ell_a$，则
\begin{equation}
K^{\rm F}=H^{\rm F}-\tau_Ql+\tau_Q\KL(Q^{\rm tail}\Vert P^{\rm tail}),
\end{equation}
无偏标量控制为 $H^{\rm F}+\tau_Qc+J_A(Q_A/\beta_A)(-\ell_A-c)$。其支持要求覆盖未知的 $Q$ 尾部。直接微分粗化正向 KL 得 $-\sum_CQ_C\bar\psi_C$，其中 $\bar\psi_C$ 仍是 $P^C$ 下的 score 均值。
\end{proof}

\section{数值估计与反传的独立核验}
\begin{proposition}[固定样本反传与积分导数]
设 $\phi$ 可微，$A\sim P_\theta$、$c_A=P_\theta(A)/\sg(P_\theta(A))$。固定样本反传满足
\begin{equation}
\nabla\phi(\ell_A)=\phi'(\ell_A)\psi_A,\qquad
\nabla[c_A\phi(\ell_A)]=[\phi(\ell_A)+\phi'(\ell_A)]\psi_A.
\end{equation}
第二式的期望为 $\nabla\E_P\phi(\ell)$。因此数值无偏与裸反传梯度无偏是不同条件；下表给出 K1--K4 与 K3$^{++}$ 的逐样本结果。
\end{proposition}
\begin{proof}
令 $A\sim P$，$\ell=\ell_A$，并定义 $c_A=P_\theta(A)/\sg(P_\theta(A))$。固定动作求导有
\begin{equation}
\nabla\phi(\ell)=\phi'(\ell)\psi,\qquad
\nabla[c_A\phi(\ell)]=[\phi(\ell)+\phi'(\ell)]\psi.
\end{equation}
第二式的期望由式~\eqref{eq:measure-derivative} 恰好恢复积分导数。数值与反传分类如下，均先讨论无截断的解析对象：
\begin{center}
\small
\begin{tabular}{lll}
\toprule
程序 & 前向值 & 固定样本反传\\
\midrule
K1 & $\ell$ & $\psi$\\
K2 & $\ell^2/2$ & $\ell\psi$\\
K3 & $e^{-\ell}+\ell-1$ & $(1-e^{-\ell})\psi$\\
K4：$c_A\sg(\ell)$ & $\ell$ & $\ell\psi$\\
live K4：$c_A\ell$ & $\ell$ & $(\ell+1)\psi$\\
K3$^{++}$：$c_A(e^{-\ell}+\ell-1)$ & $e^{-\ell}+\ell-1$ & $\ell\psi$\\
\bottomrule
\end{tabular}
\end{center}
K1、K3 的数值对逆向 KL 无偏，因为 $\E_Pe^{-\ell}=1$；K2 一般有偏，其反传却是逆向 KL 的无偏 score 梯度。K3 的裸反传期望是 $-\E_Q\psi$，对应正向 KL。K4 与 K3$^{++}$ 的逐样本反传相同，前向方差不同。Top-K 和相应 score 构造的外部背景见 EMA-PG\footfullcite{ema2026}；这里按明确公式区分命名。
\end{proof}

\begin{proposition}[K5 的双重无偏性]
完整正支持且 $A\sim P$ 时，$K5=-\sg(e^{-\ell_A})\ell_A+\log P_A-\sg(\log P_A)$ 的前向期望为 $\KL(Q\Vert P)$，固定样本反传期望为其真实参数梯度。
\end{proposition}
\begin{proof}
若需要正向 KL 的数值与 score 反传同时匹配，可取
\begin{equation}
K5=-\sg(e^{-\ell})\ell+\log P_A-\sg(\log P_A).
\end{equation}
其前向为 $-e^{-\ell}\ell$，期望为 $K^{\rm F}$；反传为 $(1-e^{-\ell})\psi$，期望为 $g^{\rm F}$。正支持和实际采样测度是这些恒等式的必要条件。
\end{proof}

\begin{proposition}[确定性头部导数与头尾常数匹配]
固定支持 $S$，令 $d_S=\nabla P_S$，则
\begin{equation}
\nabla H=h_g+d_S,\quad \nabla K_{\rm TA}=h_g-ld_S,\quad
\nabla K_{\rm norm}=h_g/P_S-Hd_S/P_S^2.
\end{equation}
对任意固定 $\upsilon$，$h_g+\upsilon d_S+J_A(\ell_A+\upsilon)\psi_A$ 无偏；只在头部或尾部保留该常数会留下对应的 $\pm d_S$ 偏差。
\end{proposition}
\begin{proof}
确定性头部数值的真实导数也应保留概率权重项。令 $d_S=\nabla P_S$，则
\begin{equation}
\nabla H=h_g+d_S,\quad
\nabla(H+\tau_Pl)=h_g-ld_S,\quad
\nabla K_{\rm norm}=\frac{h_g}{P_S}-\frac{H}{P_S^2}d_S.
\end{equation}
最后一式对商 $H/P_S$ 求导并与 $-\log P_S$ 的导数抵消得到，分母明确为标量 $P_S^2$。

头尾 score 常数必须匹配：$h_g+\upsilon d_S+J_A(\ell_A+\upsilon)\psi_A$ 对任意固定 $\upsilon$ 无偏，因为 $\E_PJ_A\psi_A=-d_S$。只在头部或尾部保留一个常数项，会分别留下 $\pm d_S$ 偏差。
\end{proof}

\begin{proposition}[固定样本势与非保守更新场]
冻结参考、头部和抽样动作。HETS 与 ResOPD 的场分别是 $H-P_S+J_A\ell_A^2/2$ 和 $K_{\rm TA}+J_A(\ell_A^2-l^2)/2$ 的普通梯度。SC-TA 的场一般非保守，其交叉偏导差为
\begin{equation}
-J_A[(\partial_i l)\psi_j-(\partial_j l)\psi_i].
\end{equation}
逐样本的保守性与总体梯度无偏性不相互替代。
\end{proposition}
\begin{proof}
可以用 $\sg(F)+L-\sg(L)$ 让前向显示 $F$、反向使用 $L$ 的导数；这不意味着前向的普通导数就是更新场。HETS 与 ResOPD 的固定样本势分别为
\begin{equation}
H-P_S+\tfrac12J_A\ell_A^2,\qquad
K_{\rm TA}+\tfrac12J_A(\ell_A^2-l^2).
\end{equation}
SC-TA 的场 $\nabla K_{\rm TA}+J_A(\ell_A-l)\nabla\ell_A$ 一般非保守，其交叉偏导差是
\begin{equation}
-J_A[(\partial_i l)\psi_j-(\partial_j l)\psi_i].
\end{equation}
取三动作 softmax、$S=\{1\},A=2$，$(i,j)=(1,2)$ 时该差为 $P_1P_3/(1-P_1)>0$。总体期望仍可为真实梯度，因为参数依赖采样测度的导数补上了额外项。
\end{proof}

\section{配分函数估计与可控展开}
\begin{proposition}[对数配分函数的 plug-in 偏差]
若 $\widehat Z>0$、$\E\widehat Z=Z>0$、$\E|\log\widehat Z|<\infty$，则 $\E\log\widehat Z\le\log Z$，非退化时严格。这个结论不排除使用其他随机构造无偏估计 $\log Z$。
\end{proposition}
\begin{proof}
固定 $q,f$，$Z=\E_qe^f$。精确头部加尾部样本均值无偏估计 $Z$，但若 $\widehat Z>0$、$\E|\log\widehat Z|<\infty$，Jensen 给
\begin{equation}
\E\log\widehat Z\le\log Z,
\end{equation}
非退化时严格。这个结论只针对 plug-in，不排除其他信息条件下无偏估计 $\log Z$。例如有限词表完全覆盖事件 $E$ 的概率 $c>0$ 已知时，$\mathbf1_E\log Z/c$ 是无偏构造，方差为 $(c^{-1}-1)(\log Z)^2$，但通常成本很高。
\end{proof}

\begin{theorem}[log-of-mean 的偏差与方差展开]
固定有限正分布与头部，$\widehat Z_M$ 为有界独立尾部变量均值加确定头部，具有与 $M$ 无关的正下界。若尾部单样本方差为 $V$，则
\begin{align}
\E\log\widehat Z_M-\log Z&=-V/(2MZ^2)+O(M^{-3/2}),\\
\Var\log\widehat Z_M&=V/(MZ^2)+O(M^{-3/2}).
\end{align}
\end{theorem}
\begin{proof}
对固定有限正分布，设 $\widehat Z_M$ 是有界独立尾部变量的均值加确定头部，且有与 $M$ 无关的正下界。令 $\epsilon=\widehat Z_M-Z$，尾部单样本方差为 $V$。则
\begin{equation}
\E\epsilon=0,\quad \E\epsilon^2=V/M,\quad
\E\epsilon^4=O(M^{-2}),\quad \E|\epsilon|^3=O(M^{-3/2}).
\end{equation}
第四矩展开仅留下四索引相同和两两配对项，分别为 $O(M)$、$O(M^2)$，除以 $M^4$ 得结论；第三绝对矩由 Cauchy--Schwarz 控制。$\log$ 的三阶导数在统一正区间有界，因此全区间 Taylor 余项满足 $|R_3|\le C|\epsilon|^3$。取期望得到
\begin{equation}
\E\log\widehat Z_M-\log Z=-\frac{V}{2MZ^2}+O(M^{-3/2}).
\end{equation}
方差中线性项贡献 $V/(MZ^2)$，其与二次项的交叉由第三矩控制，二次项平方由第四矩控制，故 $\Var\log\widehat Z_M=V/(MZ^2)+O(M^{-3/2})$。
\end{proof}

\begin{theorem}[累积量截断的可控偏差]
固定有限正分布，令 $\varphi(t)=\log\E_qe^{tr}$。用 $M\ge L$ 个独立样本无偏估计前 $L$ 阶累积量，再构造 $\widehat L=\sum_{j=1}^L\eta^j\widehat\kappa_j/j!$，则
\begin{equation}
\E\widehat L-\log Z=-\frac{\eta^{L+1}}{L!}\int_0^1(1-u)^L\varphi^{(L+1)}(u\eta)\,du.
\end{equation}
固定阶数下方差为 $O(M^{-1})$。有限阶公式不要求无限 Taylor 级数收敛。
\end{theorem}
\begin{proof}
几何信号 $f=\eta r$ 时可改用累积量展开。令 $\varphi(t)=\log\E_qe^{tr}$，指数倾斜分布为 $q_t\propto qe^{tr}$。逐项微分得 $\varphi^{(j)}(t)$ 为 $q_t$ 下第 $j$ 阶累积量。有限阶 Taylor 公式给出
\begin{equation}
\varphi(\eta)=\sum_{j=1}^L\frac{\eta^j}{j!}\varphi^{(j)}(0)
+\frac{\eta^{L+1}}{L!}\int_0^1(1-u)^L\varphi^{(L+1)}(u\eta)\,du.
\end{equation}
每阶累积量可按集合分拆 $\pi$ 写成
\begin{equation}
\kappa_j(r)=\sum_{\pi\in\mathcal P_j}(|\pi|-1)!(-1)^{|\pi|-1}
\prod_{B\in\pi}\E r^{|B|}.
\end{equation}
用互异样本索引的乘积平均替代每个矩乘积，即得到无偏估计，要求样本数 $M\ge j$。独立性逐项保证均值正确。固定阶有界核的两项乘积只在索引交叠时可能有协方差，交叠概率为 $O(1/M)$，故方差及各阶协方差均为 $O(1/M)$。截断估计的偏差正是上式积分余项的负值；无限级数还需位于收敛半径内。
\end{proof}

\begin{proposition}[头部配分函数的尾部级数估计]
若确定头部 $Z_H>0$、尾部贡献 $\tau=\E_qu$ 满足 $\tau/Z_H<1$，则以互异样本乘积无偏估计 $\tau^j$ 所得 $L$ 阶 log 尾部展开，仅保留确定截断偏差，其绝对值不超过 $(\tau/Z_H)^{L+1}/(L+1)$。样本数需 $M\ge L$。
\end{proposition}
\begin{proof}
另一种展开从确定头部 $Z_H>0$ 开始。令 $\tau=\E_qu$ 为尾部贡献，若 $\tau/Z_H<1$，则
\begin{equation}
\log(Z_H+\tau)=\log Z_H+\sum_{j=1}^L\frac{(-1)^{j+1}}j(\tau/Z_H)^j+R_{L+1},
\quad |R_{L+1}|\le\frac{(\tau/Z_H)^{L+1}}{L+1}.
\end{equation}
用互异索引乘积的 U-统计量无偏估计 $\tau^j$，样本数 $M\ge L$，留下可控的确定截断偏差。尾部均值小不自动提供逐样本小上界，实际方差仍要由核的矩计算。
\end{proof}

\section{截断、归一化与尾部匹配的差别}
\begin{proposition}[截断与 Tail-aware 的平方误差比较]
\label{prop:truncation-mse}
固定正分布 $P,Q$ 和非空真子集 $S$，令
$H=\sum_{a\in S}P_a\log(P_a/Q_a)$，
$\tau_P=P(S^c)$、$\tau_Q=Q(S^c)$，
$M=\tau_P\log(\tau_P/\tau_Q)$、
$B=\tau_P\KL(P^{S^c}\Vert Q^{S^c})\ge0$。
则完整 KL 为 $K=H+M+B$，截断值 $H$ 与 Tail-aware 值 $H+M$ 满足
\begin{equation}
(H-K)^2-(H+M-K)^2=M(M+2B).
\end{equation}
因此 Tail-aware 的下界性质不保证其平方误差小于截断值。
\end{proposition}
\begin{proof}
在尾部逐点分解
$\log(P_a/Q_a)=\log(\tau_P/\tau_Q)+
\log(P_a^{S^c}/Q_a^{S^c})$，求和得 $K=H+M+B$。
两个误差分别为 $-(M+B)$ 和 $-B$，平方相减即得公式。
$M$ 可正可负。例如取 $P=(0.4,0.2,0.4)$、
$Q=(0.2,0.4,0.4)$、$S=\{1\}$，
有 $M=0.6\log(3/4)<0$、
$B=0.2\log(2/3)+0.4\log(4/3)>0$，
且 $M+2B<0$，此时差值为正。
再保持头尾质量不变并令一个 $Q$ 尾概率趋零，$B$ 无界增长，
最终 $M+2B>0$，差值转负。两种排序均可发生。
\end{proof}

\begin{proposition}[归一化头部 KL 对尾 logit 的不敏感性]
\label{prop:tail-logit}
令 $P=\operatorname{softmax}(z)$，$Q,S$ 固定，
$P_S=\sum_{a\in S}P_a$。
归一化头部 KL $K_{\rm norm}=\KL(P^S\Vert Q^S)$ 与截断值 $H$ 满足
\begin{equation}
\frac{\partial K_{\rm norm}}{\partial z_j}=0,\qquad
\frac{\partial H}{\partial z_j}=-P_j(H+P_S),\qquad j\notin S.
\end{equation}
\end{proposition}
\begin{proof}
$P_a^S=e^{z_a}/\sum_{i\in S}e^{z_i}$ 不含尾 logit，故第一式成立。
对 $a\in S,j\notin S$，
$\partial P_a/\partial z_j=-P_aP_j$、
$\partial\log(P_a/Q_a)/\partial z_j=-P_j$。
对 $H$ 使用乘积法则并求和得第二式。
归一化会消去头部总质量，而截断求和仍通过 softmax 分母依赖尾部。
\end{proof}

\begin{proposition}[尾形状匹配时的残差梯度噪声]
\label{prop:tailmatch-noise}
设 $P^{S^c}=Q^{S^c}$，$\ell=\log(\tau_P/\tau_Q)$，
$\psi_a=\nabla\log P_a$、
$\bar\psi_C=\E_{P^{S^c}}\psi$。
前文标量桶残差修正（取 $c=\ell$）与 SC-TA 的尾形状残差在此点逐样本为零。
ResOPD 的零均值残差为
$J_A\ell(\psi_A-\bar\psi_C)$，$A\sim P$，
其协方差为
\begin{equation}
\tau_P\ell^2\Cov_{P^{S^c}}(\psi).
\end{equation}
该量可严格为正，即使 Tail-aware 的值和梯度已精确。
\end{proposition}
\begin{proof}
尾形状匹配给尾 log-ratio 恒为 $\ell$，所以
标量桶残差修正的值残差及 SC-TA 的形状残差因子均为零。
ResOPD 保留 $\ell$ 乘中心化尾 score；其尾条件均值为零。
对二阶矩按 $J_A=\mathbf1\{A\notin S\}$ 条件化，即得协方差。
例取 $P=(0.4,0.2,0.4)$、$Q=(0.7,0.1,0.2)$、
$S=\{1\}$。两者尾条件分布均为 $(1/3,2/3)$，
$\tau_P=3/5$、$\ell=\log2$。分类 softmax 的
$\psi_a=e_a-P$ 给尾协方差迹 $1-(1/3)^2-(2/3)^2=4/9$，
故残差协方差迹为 $4(\log2)^2/15>0$。
\end{proof}

\section{异分布采样与自动微分语义}
\begin{proposition}[异分布下的双控制项一般耦合]
\label{prop:offpolicy-control}
令 $A\sim\beta$，$\beta$ 为正分布，$w=P/\beta$，
$J_a=\mathbf1\{a\notin S\}$。取尾向量 $\chi_a$ 满足
$\sum_{a\notin S}P_a\chi_a=0$，定义零均值控制项
$U_{0,j}=(J_Aw_A-\tau_P)e_j$、$U_1=J_Aw_A\chi_A$。
则其 Gram 交叉项为
\begin{equation}
\E_\beta\langle U_{0,j},U_1\rangle
=\sum_{a\notin S}\frac{P_a^2}{\beta_a}\langle e_j,\chi_a\rangle.
\end{equation}
在 $\beta=P$ 时该项为零，一般 $\beta$ 下不为零。
\end{proposition}
\begin{proof}
$\E_\beta J_Aw_A=\tau_P$、
$\E_\beta J_Aw_A\chi_A=0$ 保证两个控制项均值为零。
展开内积，常数 $-\tau_P$ 对应项因后一等式消去，
其余项中的 $\beta_aw_a^2=P_a^2/\beta_a$ 给所述结果。
例如在两点全尾集合上，取 $P=(1/2,1/2)$、
$\beta=(1/4,3/4)$ 和一维 $\chi=(1,-1)$，
交叉项为 $1-1/3=2/3$。因此两个系数的联合最优值
应解完整 Gram 线性系统，不能分别按正交控制项优化。
\end{proof}

\begin{proposition}[停止梯度的双参数表示]
\label{prop:stopgrad-diagonal}
设样本 $a$ 固定，$S(\theta,\bar\theta,a)$ 对两参数可微。
冻结第二参数后在对角线取值的反传场定义为
\begin{equation}
B(\theta,a)=\partial_1S(\theta,\bar\theta,a)|_{\bar\theta=\theta}.
\end{equation}
而前向数值函数 $s(\theta,a)=S(\theta,\theta,a)$ 的普通导数为
\begin{equation}
\nabla s(\theta,a)=B(\theta,a)
+\partial_2S(\theta,\bar\theta,a)|_{\bar\theta=\theta}.
\end{equation}
若动作分布 $P_\theta$ 也随参数变化，则
\begin{equation}
\nabla\E_{P_\theta}s
=\E_{P_\theta}[\nabla s+s\nabla\log P_\theta].
\end{equation}
\end{proposition}
\begin{proof}
第一式是复合映射 $\theta\mapsto(\theta,\theta)$ 的链式法则。
第二式在有限词表上直接对 $\sum_aP_\theta(a)s(\theta,a)$ 使用乘积法则。
这两个额外项分别来自被停止的参数路径和采样测度。
因此前向数值相同、固定样本反传相同、总体梯度相同是三个需要分别核对的命题。
即使 $B$ 连续，也需满足相应可积条件才能将它称为某个标量目标的梯度。
\end{proof}

\begin{proposition}[前缀混合下的误差与方差口径]
\label{prop:prefix-mse}
设前缀 $s\sim d$，估计器 $\widehat K$ 的条件偏差为
$c(s)=\E[\widehat K\mid s]-K(s)$，条件方差为 $V(s)$。
只要二阶矩有限，
\begin{align}
\E[\widehat K-K(s)]&=\E_dc(s),\\
\E[(\widehat K-K(s))^2]&=\E_d[V(s)+c(s)^2],\\
\Var(\widehat K)&=\E_dV(s)+\Var_d(K(s)+c(s)).
\end{align}
\end{proposition}
\begin{proof}
条件化后误差均值为 $c(s)$，其平方均值为 $V(s)+c(s)^2$，
对 $d$ 平均得前两式。对估计器本身使用全方差公式得第三式。
跨前缀原始数值方差包含真实目标 $K(s)$ 的变化，
不能用它替代估计误差的均方值或条件采样方差。
\end{proof}

